% Chapter 8
\section[Beyond Straight Lines: Features, Kernels, and Foundations]{Beyond Straight Lines: Features, Kernels, and the Foundation of Everything Else}\label{sec:beyond}

We have kept one promise silently: that the model be linear in $\bbeta$. This chapter cashes in that promise. It shows (i) how ``linear'' quietly becomes ``arbitrarily flexible'' through feature engineering, (ii) how the kernel trick pushes this to infinity, and (iii) why linear regression is the load-bearing element beneath most of the models that displaced it in the headlines.

\subsection{Nonlinear relationships, linear machinery}

Recall the remark in Chapter~\ref{sec:intro}: linearity is required in the \emph{parameters}, not the raw inputs. Given any functions $\phi_1, \dots, \phi_m$ of the raw features — polynomials $\phi_j(x) = x^j$, splines, sinusoids, interactions $x_1 x_2$, indicator functions — define the \emph{basis-expanded design matrix} $\bm{\Phi}$ with rows $(\phi_1(\x_i), \dots, \phi_m(\x_i))$ and fit
\[
  y_i = \sum_{j=1}^m \beta_j \phi_j(\x_i) + \varepsilon_i
  \quad\Longleftrightarrow\quad
  \bhat = (\bm{\Phi}^\top \bm{\Phi})^{-1} \bm{\Phi}^\top \y .
\]
\emph{Every theorem of Chapters \ref{sec:ls}--\ref{sec:probability} applies verbatim}, with $\bm{\Phi}$ in place of $\X$: OLS is BLUE, the $t$- and $F$-tables are exact under Gaussian errors, and ridge regularizes the expanded basis. Taylor's theorem guarantees that smooth functions are locally well approximated by polynomials; splines piece that locality together. ``Why linear regression'' is not ``why draw straight lines'' — it is ``why solve one exactly-solvable convex problem after a wise change of coordinates.''

The one new danger is the bias--variance curve of Chapter~\ref{sec:bias-variance} tilting: with $m$ large relative to $n$, $\bm{\Phi}^\top\bm{\Phi}$ becomes ill-conditioned and variance explodes — which is precisely when ridge (or lasso) stops being optional.

\subsection{The kernel view: infinite features, finite computation}

Here is one of the most beautiful facts in machine learning. Every appearance of the data in the normal equations is through the Gram matrix $\bm{\Phi}\bm{\Phi}^\top$ (or $\bm{\Phi}^\top\bm{\Phi}$). Its entries are inner products:
\[
  [\bm{\Phi}\bm{\Phi}^\top]_{ij} = \sum_{k=1}^m \phi_k(\x_i)\, \phi_k(\x_j) =: K(\x_i, \x_j),
\]
the \emph{kernel} of the basis. For many bases — polynomials of all degrees, radial Gaussians — the function $K(\x, \x')$ has a closed form even though the basis itself is infinite-dimensional. \textbf{Kernel ridge regression} writes the solution entirely in terms of $K$:
\[
  \hat{f}(\x) = \bm{k}(\x)^\top (\bm{K} + \lambda \bm{I})^{-1} \y,
  \qquad
  K_{ij} = K(\x_i, \x_j), \;\; \bm{k}(\x) = (K(\x, \x_1), \dots, K(\x, \x_n))^\top,
\]
an $n \times n$ linear system regardless of the (possibly infinite) number of features. Linear regression on an infinite-dimensional basis is not a slogan; it is a matrix inversion. Gaussian-process regression is exactly this object dressed in probabilistic language. The lesson generalizes: \emph{methods are often separable from representations}, and the linear algebra travels better than the features.

\subsection{Logistic regression: the same geometry, a different noise model}

For binary outcomes $y_i \in \{0, 1\}$ a line predicting raw probabilities is wrong — predictions can leave $[0,1]$. The \textbf{generalized linear model} (GLM) repairs this with a link function:
\[
  \Pr(y_i = 1 \mid \x_i) = g^{-1}(\x_i^\top \bbeta), \qquad g(p) = \log \frac{p}{1-p} \;(\text{logit link}).
\]
Fitting is no longer one matrix inversion — the likelihood is not quadratic — but the \emph{iteratively reweighted least squares} (IRLS) algorithm reveals the lineage: each iteration linearizes the problem around the current fit and solves a \emph{weighted least squares} step (Chapter~\ref{sec:gauss-markov}'s GLS). Logistic regression \emph{is} linear regression, iterated, with weights and a squashing nonlinearity. The geometry (projection onto a tangent space), the theory (asymptotic normality of the MLE), and the diagnostics all inherit from the linear model.

\subsection{Neural networks: stacked, re-fit, regularized regressions}

A one-hidden-layer network with linear output is
\[
  f(\x) = \sum_{k=1}^h w_k \, \sigma(\bm{a}_k^\top \x + b_k) + c,
\]
where $\sigma$ is a nonlinearity (ReLU, sigmoid). Each hidden unit computes a \emph{learned feature} $\sigma(\bm{a}_k^\top \x + b_k)$ — basis functions $\phi_k$ whose shapes are themselves fit from data, instead of chosen by the analyst — and the output layer is literally a linear regression on those features. Training alternates: fix features, solve (regularized) linear least squares; then adjust features by gradient descent. Deep learning, from this vantage, is \emph{representation learning glued to linear regression}, with all the old themes intact: the loss is still (often) squared error, overfitting is still the bias--variance trade-off, and the output layer of every regression network is the formula \eqref{eq:ridge-sol} wearing a GPU.

\subsection{Why this chapter matters for practice}

The pragmatic upshot for someone building models today:

\begin{itemize}
  \item \textbf{Start linear.} Fit OLS on sensibly engineered features first. It is instant, interpretable, honest about its uncertainty, and provides the baseline every fancier model must beat — if a gradient-boosted ensemble cannot beat a well-specified linear model on validation data, ship the linear model.
  \item \textbf{Linearity is a coordinate choice.} Before blaming the model for underfitting, blame the coordinates: logs, interactions, splines, and domain transforms move most ``nonlinear'' problems back into the exactly-solvable regime.
  \item \textbf{The theorems travel.} BLUE, the $t$-table, shrinkage, and cross-validation recur, appropriately modified, in nearly every classical model. Learning them once, properly, in the linear case is the highest-leverage hour in applied statistics.
\end{itemize}

\begin{keybox}[Chapter verdict]
``Linear'' never meant ``simple.'' Linear regression plus features equals nonlinear regression with exact solutions and honest inference; plus kernels equals infinite flexibility in $n \times n$ algebra; iterated and squashed equals logistic regression; stacked and gradient-trained equals neural networks. The most modern models are best understood as answers to the question this tutorial started with: why linear regression?
\end{keybox}
